% Options for packages loaded elsewhere
% Options for packages loaded elsewhere
\PassOptionsToPackage{unicode}{hyperref}
\PassOptionsToPackage{hyphens}{url}
\PassOptionsToPackage{dvipsnames,svgnames,x11names}{xcolor}
%
\documentclass[
  german,
  10pt,
  twoside]{article}
\usepackage{xcolor}
\usepackage[top=2.5cm,bottom=2.5cm,left=2cm,right=2cm]{geometry}
\usepackage{amsmath,amssymb}
\setcounter{secnumdepth}{5}
\usepackage{iftex}
\ifPDFTeX
  \usepackage[T1]{fontenc}
  \usepackage[utf8]{inputenc}
  \usepackage{textcomp} % provide euro and other symbols
\else % if luatex or xetex
  \usepackage{unicode-math} % this also loads fontspec
  \defaultfontfeatures{Scale=MatchLowercase}
  \defaultfontfeatures[\rmfamily]{Ligatures=TeX,Scale=1}
\fi
\usepackage{lmodern}
\ifPDFTeX\else
  % xetex/luatex font selection
\fi
% Use upquote if available, for straight quotes in verbatim environments
\IfFileExists{upquote.sty}{\usepackage{upquote}}{}
\IfFileExists{microtype.sty}{% use microtype if available
  \usepackage[]{microtype}
  \UseMicrotypeSet[protrusion]{basicmath} % disable protrusion for tt fonts
}{}
\usepackage{setspace}
\makeatletter
\@ifundefined{KOMAClassName}{% if non-KOMA class
  \IfFileExists{parskip.sty}{%
    \usepackage{parskip}
  }{% else
    \setlength{\parindent}{0pt}
    \setlength{\parskip}{6pt plus 2pt minus 1pt}}
}{% if KOMA class
  \KOMAoptions{parskip=half}}
\makeatother
% Make \paragraph and \subparagraph free-standing
\makeatletter
\ifx\paragraph\undefined\else
  \let\oldparagraph\paragraph
  \renewcommand{\paragraph}{
    \@ifstar
      \xxxParagraphStar
      \xxxParagraphNoStar
  }
  \newcommand{\xxxParagraphStar}[1]{\oldparagraph*{#1}\mbox{}}
  \newcommand{\xxxParagraphNoStar}[1]{\oldparagraph{#1}\mbox{}}
\fi
\ifx\subparagraph\undefined\else
  \let\oldsubparagraph\subparagraph
  \renewcommand{\subparagraph}{
    \@ifstar
      \xxxSubParagraphStar
      \xxxSubParagraphNoStar
  }
  \newcommand{\xxxSubParagraphStar}[1]{\oldsubparagraph*{#1}\mbox{}}
  \newcommand{\xxxSubParagraphNoStar}[1]{\oldsubparagraph{#1}\mbox{}}
\fi
\makeatother


\usepackage{graphicx}
\makeatletter
\newsavebox\pandoc@box
\newcommand*\pandocbounded[1]{% scales image to fit in text height/width
  \sbox\pandoc@box{#1}%
  \Gscale@div\@tempa{\textheight}{\dimexpr\ht\pandoc@box+\dp\pandoc@box\relax}%
  \Gscale@div\@tempb{\linewidth}{\wd\pandoc@box}%
  \ifdim\@tempb\p@<\@tempa\p@\let\@tempa\@tempb\fi% select the smaller of both
  \ifdim\@tempa\p@<\p@\scalebox{\@tempa}{\usebox\pandoc@box}%
  \else\usebox{\pandoc@box}%
  \fi%
}
% Set default figure placement to htbp
\def\fps@figure{htbp}
\makeatother



\ifLuaTeX
\usepackage[bidi=basic]{babel}
\else
\usepackage[bidi=default]{babel}
\fi
% get rid of language-specific shorthands (see #6817):
\let\LanguageShortHands\languageshorthands
\def\languageshorthands#1{}
\ifLuaTeX
  \usepackage[german]{selnolig} % disable illegal ligatures
\fi


\setlength{\emergencystretch}{3em} % prevent overfull lines

\providecommand{\tightlist}{%
  \setlength{\itemsep}{0pt}\setlength{\parskip}{0pt}}



 
\usepackage[style=authoryear,backend=biber,style=authoryear-comp,uniquename=false,uniquelist=false]{biblatex}
\addbibresource{../references.bib}


\usepackage{csquotes}
\usepackage{booktabs}
\usepackage{array}
\usepackage{multirow}
\usepackage{wrapfig}
\usepackage{float}
\usepackage{colortbl}
\usepackage{pdflscape}
\usepackage{threeparttable}
\usepackage{threeparttablex}
\usepackage[normalem]{ulem}
\usepackage{makecell}
\usepackage{xcolor}
\usepackage{multicol}
\usepackage{fancyhdr}
\usepackage{titlesec}
\usepackage{etoolbox}
\usepackage{tabularx}
\usepackage{afterpage}
\usepackage{tikz}
\usepackage{longtable}



% Suppress default maketitle
\renewcommand{\maketitle}{}

% Define WISTA colors (Havelock Blue as main)
\definecolor{wistaBlue}{RGB}{65,125,180}
\definecolor{wistaLightBlue}{RGB}{120,170,215}
\definecolor{wistaDarkGray}{RGB}{64,64,64}
\definecolor{wistaLightGray}{RGB}{128,128,128}

% Section formatting (##) - number 20% bigger, lines closer, text closer to bottom line
% e.g., "2" with lines around "Einleitung"
\titleformat{\section}[display]
  {\normalfont\large\bfseries}
  {\color{wistaBlue}\Large\bfseries\thesection\\\noindent\rule{\columnwidth}{1.5pt}}
  {0pt}
  {\vspace{0pt}\color{wistaDarkGray}\large\bfseries}
  [\vspace{0pt}\noindent{\color{wistaBlue}\rule{\columnwidth}{1.5pt}}\vspace{0.05cm}]
\titlespacing*{\section}{0pt}{1.5ex plus 0.3ex minus .1ex}{0.2ex plus .1ex}

% Subsection formatting (###) - grey text with blue line under
% e.g., "2.1 Handlungsbedarf"
\renewcommand{\thesubsection}{\thesection.\arabic{subsection}}
\titleformat{\subsection}
  {\color{wistaDarkGray}\normalsize\bfseries}
  {\color{wistaDarkGray}\normalsize\bfseries\thesubsection}
  {0.5em}
  {\color{wistaDarkGray}\normalsize\bfseries}
  [\vspace{0pt}\noindent{\color{wistaBlue}\rule{\columnwidth}{0.8pt}}\vspace{0.05cm}]
\titlespacing*{\subsection}{0pt}{1ex plus 0.2ex minus .1ex}{0.2ex plus .1ex}
  
% Subsubsection formatting (####) - grey text, no underline
% e.g., "2.1.1 Steigende Datenmenge"
\renewcommand{\thesubsubsection}{\thesubsection.\arabic{subsubsection}}
\titleformat{\subsubsection}
  {\color{wistaDarkGray}\normalsize\bfseries}
  {\color{wistaDarkGray}\normalsize\bfseries\thesubsubsection}
  {0.5em}
  {\color{wistaDarkGray}\normalsize\bfseries}

% Paragraph formatting (4th level) - numbered as subsection.subsubsection.paragraph
\renewcommand{\theparagraph}{\thesubsubsection.\arabic{paragraph}}
\titleformat{\paragraph}[runin]
  {\color{wistaDarkGray}\normalsize\bfseries}
  {\color{wistaBlue}\normalsize\bfseries\theparagraph}
  {0.5em}
  {\color{wistaDarkGray}\normalsize\bfseries}
  [.]
\titlespacing*{\paragraph}{0pt}{1ex plus 0.5ex minus 0.2ex}{0.5em}

% Header and footer setup with blue line under header
\pagestyle{fancy}
\fancyhf{}
\fancyhead[L]{\small\textit{Oliver Hauke, Annette Kristiansen}}
\fancyhead[R]{\small\textit{Effiziente Analyse großer Forschungsdaten}}
\fancyfoot[L]{\small\thepage}
\fancyfoot[R]{\small Statistisches Bundesamt | WISTA | {\color{wistaBlue}6} | 2025}
\renewcommand{\headrulewidth}{0pt}
\renewcommand{\footrulewidth}{0pt}
% Add blue line under header
\renewcommand{\headrule}{\color{wistaBlue}\hrule width\headwidth height 1pt \vskip 2pt}
\renewcommand{\headrulewidth}{1pt}

% Make all figures span both columns (full width)
\let\origfigure\figure
\let\endorigfigure\endfigure
\renewenvironment{figure}[1][htbp]{%
  \origfigure*[#1]%
}{%
  \endorigfigure*%
}

% Allow manual control of table* environment for full-width tables
% Tables by default stay in single column, but can use table* manually



% Make bibliography full-width (one column) - for biblatex
\defbibheading{bibliography}[\bibname]{%
  \onecolumn
  \section*{#1}%
  \markboth{#1}{#1}%
}



% Footnotes stay within column
\usepackage[stable]{footmisc}

% Custom footnote formatting: "|^1" in dark WISTA gray
\renewcommand{\thefootnote}{{\color{wistaDarkGray}|$^{\arabic{footnote}}$}}

% Pifont for dingbat symbols
\usepackage{pifont}

% Better table handling for two-column mode
\usepackage{array,tabularx,calc}

% Custom WISTA table format
\usepackage{caption}
\usepackage{xcolor}
\usepackage{colortbl}
\usepackage{array}

% Define custom caption format for WISTA tables
\DeclareCaptionFormat{wista}{%
  {\color{wistaLightBlue}\bfseries #1#2}\par\vspace{0.2em}%
  {\color{wistaDarkGray}\normalfont #3}\par\vspace{0.3em}%
}

% Set up table captions with custom WISTA format
\captionsetup[table]{
  position=top,
  format=wista,
  justification=justified,
  singlelinecheck=false,
  labelsep=none,
  skip=0.3em
}

% Better table spacing and white inner lines
\setlength{\tabcolsep}{6pt}
\renewcommand{\arraystretch}{1.15}

% Set white lines for tables by default, blue for first column
\AtBeginDocument{
  \arrayrulecolor{white}
  \setlength{\arrayrulewidth}{0.3pt}
}

% Define commands for colored table borders
\newcommand{\bluevrule}{\color{wistaBlue}\vrule}
\newcommand{\whitevrule}{\color{white}\vrule}

% Define WISTA table colors
\definecolor{wistaTableBg}{RGB}{240,248,255}     % Very light blue background
\definecolor{wistaHeaderBg}{RGB}{220,230,240}    % Light blue header background
\definecolor{wistaFirstCol}{RGB}{65,125,180}     % Blue for first column lines

% Custom WISTA table command
\newcommand{\wistaTableStart}{%
  \rowcolors{2}{wistaTableBg}{wistaTableBg}%
  \arrayrulecolor{wistaFirstCol}%
}

% Custom commands for table styling
\newcommand{\wistaHeaderRow}[1]{%
  \rowcolor{wistaHeaderBg}%
  \arrayrulecolor{wistaDarkGray}%
  #1 \\
  \arrayrulecolor{white}%
}

\newcommand{\wistaFirstColRule}{\arrayrulecolor{wistaFirstCol}}
\newcommand{\wistaWhiteRule}{\arrayrulecolor{white}}

% Custom bullet points with > character - 35% width, 95% height, black extra bold, half line spacing
\usepackage{enumitem}
\setlist[itemize]{label={\scalebox{0.35}[0.95]{\fontseries{bx}\selectfont >}}, leftmargin=1.5em, itemsep=5pt, topsep=3pt, parsep=0pt}

% Redefine ref to include pifont 247 (arrow) for figure/table/section references
\let\oldref\ref
\renewcommand{\ref}[1]{\oldref{#1}\raisebox{-0.2ex}{\scriptsize\color{wistaBlue}\ding{247}}}
\makeatletter
\@ifpackageloaded{caption}{}{\usepackage{caption}}
\AtBeginDocument{%
\ifdefined\contentsname
  \renewcommand*\contentsname{Inhaltsverzeichnis}
\else
  \newcommand\contentsname{Inhaltsverzeichnis}
\fi
\ifdefined\listfigurename
  \renewcommand*\listfigurename{Abbildungsverzeichnis}
\else
  \newcommand\listfigurename{Abbildungsverzeichnis}
\fi
\ifdefined\listtablename
  \renewcommand*\listtablename{Tabellenverzeichnis}
\else
  \newcommand\listtablename{Tabellenverzeichnis}
\fi
\ifdefined\figurename
  \renewcommand*\figurename{Abbildung}
\else
  \newcommand\figurename{Abbildung}
\fi
\ifdefined\tablename
  \renewcommand*\tablename{Tabelle}
\else
  \newcommand\tablename{Tabelle}
\fi
}
\@ifpackageloaded{float}{}{\usepackage{float}}
\floatstyle{ruled}
\@ifundefined{c@chapter}{\newfloat{codelisting}{h}{lop}}{\newfloat{codelisting}{h}{lop}[chapter]}
\floatname{codelisting}{Listing}
\newcommand*\listoflistings{\listof{codelisting}{Listingverzeichnis}}
\makeatother
\makeatletter
\makeatother
\makeatletter
\@ifpackageloaded{caption}{}{\usepackage{caption}}
\@ifpackageloaded{subcaption}{}{\usepackage{subcaption}}
\makeatother
\usepackage{bookmark}
\IfFileExists{xurl.sty}{\usepackage{xurl}}{} % add URL line breaks if available
\urlstyle{same}
\hypersetup{
  pdftitle={EFFIZIENTE AUSWERTUNG DES BUSINESS-TAX-PANELS MITHILFE VON R},
  pdfauthor={Oliver Hauke; Annette Kristiansen},
  pdflang={de},
  colorlinks=true,
  linkcolor={blue},
  filecolor={Maroon},
  citecolor={Blue},
  urlcolor={Blue},
  pdfcreator={LaTeX via pandoc}}


\title{EFFIZIENTE AUSWERTUNG DES BUSINESS-TAX-PANELS MITHILFE VON R}
\author{Oliver Hauke \and Annette Kristiansen}
\date{2025-12-04}
\begin{document}
\maketitle

% Custom WISTA-style title page
\begin{titlepage}
\thispagestyle{empty}
\pagestyle{empty}
\onecolumn

% Horizontal blue line at top (header line)
\noindent{\color{wistaBlue}\rule{\textwidth}{1pt}}

\vspace{-0.5cm}

% Vertical blue line on the left (6.3cm from left edge, starts below header line, stops above footer)
\begin{tikzpicture}[remember picture, overlay]
  \draw[wistaBlue, line width=2pt] 
    ([xshift=6.3cm, yshift=-4cm]current page.north west) -- 
    ([xshift=6.3cm, yshift=3.5cm]current page.south west);
\end{tikzpicture}

\vspace*{0.5cm}

% Left side - short author bios (right-aligned to the vertical line)
\begin{minipage}[t]{3.9cm}
\selectlanguage{ngerman}%
\raggedleft
\hyphenpenalty=0
\tolerance=1000
\fontsize{7}{8.5}\selectfont
\vspace{0pt}%
{\fontsize{8}{9.5}\selectfont\color{wistaBlue}\textbf{Oliver Hauke}}\\[0.15cm]
ist Mathematiker und IT-Referent im Kompetenzzentrum „Analyse und Auswertung" des Statistischen Bundesamtes. Er verantwortet die Weiterentwicklung der technischen Infrastruktur des Forschungsdatenzentrums und berät das Netzwerk empirische Steuerforschung in der effizienten Nutzung der Systeme.\\[0.5cm]

{\fontsize{8}{9.5}\selectfont\color{wistaBlue}\textbf{Annette Kristiansen}}\\[0.15cm]
ist Ökonomin und Referentin im Referat „Unternehmenssteuern" des Statistischen Bundesamtes. Sie beschäftigt sich mit der Zusammenführung der Unternehmenssteuerstatistiken und betreut das Netzwerk empirische Steuerforschung. Für ihre externe Promotion an der Universität Bremen untersucht sie die technologische Vielfalt multinationaler Unternehmen.

\end{minipage}%
\hspace{0.4cm}% Gap to vertical line (left side)
\hspace{0.4cm}% Gap from vertical line (right side)
% Right side - title, authors, keywords, abstracts
\begin{minipage}[t]{11.8cm}
\raggedright
\vspace{1.2cm}%

% Title (uppercase via command, not bold)
{\LARGE\color{wistaDarkGray}\begin{spacing}{1.3}\MakeUppercase{Effiziente Auswertung des Business-Tax-Panels mithilfe von R}\end{spacing}}

\vspace{-0.2cm}

% Blue horizontal line (underline for title)
{\color{wistaBlue}\rule{\linewidth}{0.5pt}}

\vspace{0.5cm}

% Authors inline, dark gray
{\large\color{wistaDarkGray} Oliver Hauke, Annette Kristiansen}

\vspace{0.05cm}

% Blue line (underline for authors)
{\color{wistaBlue}\rule{\linewidth}{0.5pt}}

\vspace{0.5cm}

% Keywords (pifont 247 arrow then bold label and blue text)
\hangindent=1.2em\hangafter=1
{\color{wistaBlue}\ding{247}\ {\bfseries Schlüsselwörter:} Effiziente Analysen, R-Programmierung, Forschungsdaten, Unternehmenssteuerstatistik, Business-Tax-Panel, Netzwerk empirische Steuerforschung}

\vspace{0.3cm}

% Zusammenfassung (uppercase via command)
{\color{wistaDarkGray}\bfseries\MakeUppercase{Zusammenfassung}}\\[0.15cm]
\small
Das Forschungsdatenzentrum des Statistischen Bundesamtes baut seine technische Infrastruktur gezielt aus, um der Wissenschaft erweiterte Analysemöglichkeiten amtlicher Daten bereitzustellen. Seit November 2025 steht Forschenden exklusiver Zugang zu diesen erweiterten Systemressourcen in R und Stata zur Verfügung. R-basierte Tools -- insbesondere Apache Arrow, Parquet und DuckDB -- ermöglichen es, gängige Analysen großer Forschungsdatensätze in Echtzeit auszuführen. Am Beispiel des Business-Tax-Panel (2013 bis 2019) werden die erzielbaren Effizienzgewinne bei der Auswertung umfangreicher steuerstatistischer Daten im Rahmen des Netzwerks empirische Steuerforschung demonstriert.

\vspace{0.5cm}

% English keywords and abstract (italicized)
\hangindent=1.2em\hangafter=1
{\itshape\color{wistaBlue}\ding{247}\ {\bfseries Keywords:} Efficient analysis -- R programming -- research data -- corporate tax statistics -- Business Tax Panel -- empirical tax research network}

\vspace{0.3cm}

{\itshape\color{wistaDarkGray}\bfseries\MakeUppercase{Abstract}}\\[0.15cm]
\small\itshape
The Research Data Centre at the Federal Statistical Office is strategically expanding its technical infrastructure to provide researchers with enhanced analytical capabilities for official data. Since November 2025, researchers have had exclusive access to these extended system resources in R and Stata. R-based tools -- particularly Apache Arrow, Parquet, and DuckDB -- enable real-time analysis of large research datasets. Using the Business Tax Panel (2013--2019), this paper demonstrates the achievable efficiency gains in analysing large-scale tax statistics within the empirical tax research network.

\end{minipage}

\vfill

\vspace{0.5cm}
% Footer matching document style
\noindent\makebox[\textwidth]{%
  \small 1%
  \hfill%
  \small Statistisches Bundesamt | WISTA | {\color{wistaBlue}6} | 2025%
}

\end{titlepage}

% Stay in one-column mode for TOC (will be switched later)
\pagestyle{fancy}

\renewcommand*\contentsname{Inhaltsverzeichnis}
{
\hypersetup{linkcolor=}
\setcounter{tocdepth}{3}
\tableofcontents
}

\setstretch{1}
\section{Infrastruktur des FDZ-Bund}\label{sec-infrastruktur}

Erwähnen: Klassisch ebenso SAS- und Stata-Server

Womit wir arbeiten können\ldots{}

TODO: Vorbereitung der Daten auf neue Methoden, Werkzeuge

\subsection{StBA}\label{stba}

StBA R-Server mit noch größeren Kapazitäten und grafischen
Rechenbeschleunigern.

\subsection{FDZ OnSite Server}\label{fdz-onsite-server}

Neu seit November 2025 sind im FDZ OnSite-Server Bund verfügbar.

Zugriff zu Analyseservern in umfangreichen Kapazitäten.

IM FDZ Bund Für Stata oder R in KDFV oder GWAP stehen jeder Nutzung zur
Verfügung:

\textbf{Systemspezifikationen:}

\begin{itemize}
\tightlist
\item
  \textbf{Arbeitsspeicher:} bis zu 512 GB
\item
  \textbf{Kernanzahl:} 16\\
\item
  \textbf{CPU-Geschwindigkeit:} 3.1 GHz
\end{itemize}

\begin{itemize}
\tightlist
\item
  schnelle Direktverbindung zu Netzwerkspeicher
\item
  Fußnote RAM in R 128 GB kann erweitert werden.
\end{itemize}

Intern ist R in größeren Kapazitäten nutzbar und Stata nur lokal.

Nonetheless, the idea carries over to intermediate datasets stored in
your project directory and highlights the fact that small changes in
your code can result in large performance gains. This also shows that
common tasks can usually be accomplished much faster using R than using
Stata. However, we also pointed out that you should not only focus on
computing time alone as shorter computing times might be offset by
longer programming time

Bereitstellung von Paketen und organisatorische einordnung. Veränderung
der Infrastruktur langwierig, Beschaffung via Beschaffungsamt etc.
-\textgreater{} wir müssen für die Zukunft planen und umsetzen.

\subsection{Ausgangslage}\label{ausgangslage}

SAS

\begin{itemize}
\tightlist
\item
  verbreitet
\item
  geeignet für größere Datenmengen (langsam)
\end{itemize}

Stata

\begin{itemize}
\tightlist
\item
  beliebt in Wiss.
\item
  hat entsprechend zugeschnittene Features
\item
  Bei RAM-Grenze ist Schluss
\end{itemize}

\subsection{KI Anhang}\label{ki-anhang}

\subsubsection{Die Forschungsdatenzentren der amtlichen
Statistik}\label{die-forschungsdatenzentren-der-amtlichen-statistik}

Die Forschungsdatenzentren (FDZ) der amtlichen Statistik bestehen aus
zwei organisatorisch voneinander getrennten Einrichtungen: dem
\textbf{Forschungsdatenzentrum des Statistischen Bundesamts} und dem
\textbf{Forschungsdatenzentrum der Statistischen Ämter der Länder}.
Beide FDZ arbeiten eng zusammen und bieten ein abgestimmtes Daten- und
Dienstleistungsangebot für die wissenschaftliche Nutzung von Mikrodaten
der amtlichen Statistiken.

\textbf{Entstehung und Entwicklung.}

Die Gründung der FDZ geht auf ein Gutachten der „Kommission zur
Verbesserung der informationellen Infrastruktur zwischen Wissenschaft
und Statistik'' (KVI) von 1999 zurück. Mit finanzieller Förderung des
Bundesministeriums für Bildung und Forschung (BMBF) wurde im Herbst 2001
zunächst das FDZ des Statistischen Bundesamts gegründet. Das FDZ der
Statistischen Ämter der Länder folgte im April 2002.

\textbf{Aufgaben und Ziele.}

Die FDZ verfolgen das Ziel, den Zugang zu Mikrodaten der amtlichen
Statistik für die Wissenschaft nicht nur zu ermöglichen, sondern
fortlaufend zu verbessern und an die sich ändernden Bedarfe anzupassen:

\begin{itemize}
\tightlist
\item
  \textbf{Datenzugang:} Bereitstellung hochwertiger Mikrodaten für die
  wissenschaftliche Forschung
\item
  \textbf{Datenschutz:} Gewährleistung des Schutzes personenbezogener
  und unternehmensbezogener Daten nach § 16 BStatG
\item
  \textbf{Beratung:} Unterstützung der Forschenden bei der Datennutzung
  und methodischen Fragen
\item
  \textbf{Infrastruktur:} Bereitstellung einer sicheren und
  leistungsfähigen Auswertungsumgebung
\item
  \textbf{Regionalisierung:} Flächendeckende Standorte in ganz
  Deutschland für kurze Anfahrtswege
\end{itemize}

\subsubsection{Technische Infrastruktur}\label{technische-infrastruktur}

\textbf{Hardware-Ausstattung.}

Die technische Infrastruktur des FDZ-Bund umfasst verschiedene
Systemkomponenten:

\begin{table*}[htbp]
\caption{Technische Ausstattung des FDZ-Bund (beispielhafte Konfiguration)}
\label{tbl-fdz-infrastructure}
\centering
\begin{tabular*}{\textwidth}{@{\extracolsep{\fill}}llc@{}}
\toprule
Komponente & Spezifikation & Anzahl\\
\midrule
Rechenknoten (Standard) & 2x Intel Xeon Gold 6248R (48 Kerne), 384 GB RAM & 12\\
Rechenknoten (High-Memory) & 2x Intel Xeon Platinum 8280 (56 Kerne), 1,5 TB RAM & 4\\
Speichersystem & NAS-System mit 500 TB nutzbarem Speicher & 1\\
Netzwerk & 10 Gbit/s Ethernet, InfiniBand für parallele Operationen & -\\
Backup-System & Tägliches Backup, räumlich getrennte Archivierung & -\\
\bottomrule
\end{tabular*}
\end{table*}

\textbf{Software-Umgebung.}

Den Forschenden stehen verschiedene Software-Tools zur Verfügung:

\begin{itemize}
\tightlist
\item
  \textbf{Statistische Software:} R, Python, Stata, SAS, SPSS
\item
  \textbf{Datenbanksysteme:} PostgreSQL, DuckDB
\item
  \textbf{Entwicklungsumgebungen:} RStudio Server, JupyterLab
\item
  \textbf{Versionskontrolle:} Git (lokal, ohne externe Verbindung)
\end{itemize}

\subsubsection{Infrastrukturkonzept}\label{infrastrukturkonzept}

\textbf{Fachlich zentralisierte Datenhaltung.}

In Deutschland wird der überwiegende Teil der amtlichen
Statistikproduktion dezentral in den Statistischen Ämtern der Länder
durchgeführt (über 90\% aller amtlichen Statistiken). Da sich
wissenschaftliche Analysen in der Regel jedoch auf mehrere Bundesländer
oder das gesamte Bundesgebiet beziehen, haben die FDZ eine
\textbf{fachlich zentralisierte Datenhaltung} eingerichtet. Diese
ermöglicht es, die Daten des gesamten Bundesgebiets länderübergreifend
an jedem Standort bereitzustellen.

\textbf{Regionalisierte Infrastruktur.}

Die FDZ verfolgen das Ziel, immer in der Nähe der Wissenschaft zu sein.
Dies wird durch über ganz Deutschland verteilte Standorte erreicht.
Standorte befinden sich u.a. in:

\begin{itemize}
\tightlist
\item
  Berlin, Bonn, Wiesbaden (Bund)
\item
  Bad Ems, Bremen, Dresden, Düsseldorf, Erfurt, Frankfurt, Hamburg,
  Hannover, Kiel, München, Stuttgart u.v.m. (Länder)
\end{itemize}

\subsubsection{Zugangsformen und
Nutzungsmodelle}\label{zugangsformen-und-nutzungsmodelle}

Die FDZ bieten verschiedene Zugangswege an, über die unterschiedlich
anonymisierte Datenprodukte bereitgestellt werden:

\textbf{Off-Site-Zugang.}

\begin{itemize}
\tightlist
\item
  \textbf{Public Use Files (PUF):} Stark anonymisierte Daten für freie
  Nutzung
\item
  \textbf{Scientific Use Files (SUF):} Faktisch anonymisierte Daten für
  wissenschaftliche Zwecke
\item
  \textbf{Remote Scientific Use Files (Remote-SUF):} Neuer Zugangsweg
  mit höherer Detailtiefe
\end{itemize}

\textbf{On-Site-Zugang.}

\begin{itemize}
\tightlist
\item
  \textbf{Gastwissenschaftsarbeitsplatz (GWAP):} Zugang zu formal
  anonymisierten Daten vor Ort

  \begin{itemize}
  \tightlist
  \item
    Vorteile: Höchste Datendetailtiefe, direkte Beratung
  \item
    Anforderungen: Terminvereinbarung, Zutrittskontrolle,
    Output-Kontrolle
  \end{itemize}
\item
  \textbf{Kontrollierte Datenfernverarbeitung (KDFV):} Remote-Ausführung
  von Analyseskripten

  \begin{itemize}
  \tightlist
  \item
    Vorteile: Keine Anreise, zeitlich flexibel
  \item
    Einschränkungen: Keine interaktive Arbeit, Ergebnisfreigabe
    erforderlich
  \end{itemize}
\end{itemize}

\subsubsection{Datenschutz und
Sicherheit}\label{datenschutz-und-sicherheit}

\textbf{Mehrstufiges Sicherheitskonzept.}

Das FDZ-Bund implementiert ein umfassendes Sicherheitskonzept:

\begin{enumerate}
\def\labelenumi{\arabic{enumi}.}
\tightlist
\item
  \textbf{Physische Sicherheit:} Zutrittskontrolle, gesicherte
  Räumlichkeiten
\item
  \textbf{Netzwerksicherheit:} Isolierte Netzwerke, Firewalls, Intrusion
  Detection
\item
  \textbf{Zugriffsrechte:} Rollenbasierte Zugriffskontrolle,
  Authentifizierung
\item
  \textbf{Output-Kontrolle:} Prüfung aller Ergebnisse vor Freigabe
\item
  \textbf{Audit-Logging:} Protokollierung aller Zugriffe und Operationen
\end{enumerate}

\textbf{Output-Kontrolle.}

Bevor Forschungsergebnisse das FDZ verlassen dürfen, durchlaufen sie
eine mehrstufige Kontrolle:

\begin{itemize}
\tightlist
\item
  \textbf{Automatische Prüfung:} Screening auf kritische Werte (z.B.
  Fallzahlen \textless{} 5)
\item
  \textbf{Manuelle Prüfung:} Sichtung durch geschulte Mitarbeiter
\item
  \textbf{Dokumentation:} Nachvollziehbare Begründung der
  Freigabeentscheidung
\end{itemize}

\subsubsection{Performanz-Optimierung in der
FDZ-Umgebung}\label{performanz-optimierung-in-der-fdz-umgebung}

\textbf{Herausforderungen.}

Die Arbeit in der geschützten FDZ-Umgebung stellt besondere
Anforderungen:

\begin{itemize}
\tightlist
\item
  \textbf{Begrenzte Ressourcen:} Mehrere Nutzer teilen sich die
  verfügbare Hardware
\item
  \textbf{Keine Cloud-Dienste:} Externe Rechenressourcen können nicht
  genutzt werden
\item
  \textbf{Datengröße:} Große Datensätze müssen effizient verarbeitet
  werden
\item
  \textbf{Reproduzierbarkeit:} Analysen müssen dokumentiert und
  nachvollziehbar sein
\end{itemize}

\textbf{Optimierungsstrategien.}

Verschiedene Strategien helfen, die Performanz zu optimieren:

\begin{figure}

\centering{

\pandocbounded{\includegraphics[keepaspectratio]{03-infrastruktur_files/figure-pdf/fig-optimization-strategies-1.pdf}}

}

\caption{\label{fig-optimization-strategies}Strategien zur
Performanz-Optimierung im FDZ-Bund}

\end{figure}%

\subsubsection{Akkreditierung und
Qualitätssicherung}\label{akkreditierung-und-qualituxe4tssicherung}

Die Forschungsdatenzentren der Statistischen Ämter des Bundes und der
Länder sind durch den \textbf{Rat für Sozial- und WirtschaftsDaten
(RatSWD) akkreditiert}. Der RatSWD akkreditiert seit 2008
Forschungsdatenzentren, die den Zugang zu qualitativ hochwertigen
sensiblen Daten für die Wissenschaft ermöglichen und dabei die
erforderlichen Kriterien erfüllen.

\subsubsection{Best Practices für
FDZ-Nutzende}\label{best-practices-fuxfcr-fdz-nutzende}

Aus der Erfahrung mit den FDZ lassen sich folgende Empfehlungen für
effizientes Arbeiten ableiten:

\begin{enumerate}
\def\labelenumi{\arabic{enumi}.}
\tightlist
\item
  \textbf{Datenreduktion:} Frühzeitige Filterung auf relevante
  Beobachtungen und Variablen
\item
  \textbf{Inkrementelle Entwicklung:} Erst mit Stichproben arbeiten,
  dann auf Vollsatz skalieren
\item
  \textbf{Effiziente Datenformate:} Nutzung moderner Formate
  (Arrow/Parquet) statt CSV
\item
  \textbf{Dokumentation:} Sorgfältige Dokumentation des Analyseprozesses
  für Output-Kontrolle
\item
  \textbf{Ressourcenplanung:} Bewusste Planung rechenintensiver
  Operationen bei geteilter Infrastruktur
\item
  \textbf{Beratung nutzen:} Frühzeitige Kontaktaufnahme mit
  FDZ-Mitarbeitern bei methodischen Fragen
\end{enumerate}


\printbibliography



\end{document}
